% ===========================================================
%  Chapter 9 — Actor-Critic Methods and Proximal Policy Optimization
% ===========================================================
\chapter[Методы актор-критик и PPO]{Методы актор-критик и Proximal Policy Optimization}
\label{ch:actor-critic}

\begin{quote}
\itshape
Идея наличия двух взаимодействующих компонентов\,---\,одного для поведения и одного для критики\,---\,кажется столь же старой, как и сама идея обучения с подкреплением.
\upshape\hfill---Ричард С.\ Саттон и Эндрю Дж.\ Барто, \textit{Reinforcement Learning: An Introduction}
\end{quote}

\bigskip

В главе~\ref{ch:pg} была разработана теорема о градиенте стратегии и алгоритм REINFORCE, установившие, что можно напрямую оптимизировать параметризованную стратегию, следуя градиенту ожидаемой отдачи. REINFORCE концептуально элегантен, но страдает от фундаментальной статистической слабости: его оценщик градиента использует полную дисконтированную отдачу $G_t$ в качестве сигнала для каждого совершённого в эпизоде действия. Поскольку $G_t$ является суммой многих случайных вознаграждений, простирающейся до конца эпизода, его дисперсия растёт с горизонтом, что приводит к медленному и зашумлённому обучению.

Предлагаемое в данной главе средство\,---\,заменить $G_t$ оценкой с меньшей дисперсией того, насколько конкретное действие было лучше среднего\,---\,\emph{функцией преимущества}\,---\,и обучать эту оценку онлайн с помощью второй нейронной сети, называемой \emph{критиком}. Возникающий класс алгоритмов известен как методы \textbf{актор-критик}.
\index{actor-critic}
Актор (стратегия) и критик (функция ценности) взаимодействуют на каждом шаге: критик оценивает выборы актора, а актор использует эту оценку для обновления своих параметров.

На этой основе мы строим всё более мощные алгоритмы. Обобщённая оценка преимущества (GAE) обеспечивает гибкий компромисс между смещением и дисперсией для сигнала преимущества. Асинхронный актор-критик с преимуществом (A3C) использует параллельных работников для декорреляции опыта и ускорения обучения. Оптимизация стратегии с доверительной областью (TRPO) вводит ключевое понятие о том, что обновления стратегии должны оставаться близкими к текущей стратегии. Proximal Policy Optimization (PPO) достигает того же эффекта с помощью простого механизма отсечения, который работает быстро, легко реализуется и оказывается замечательно устойчивым на практике. Наконец, мы изучаем, как PPO применяется для дообучения больших языковых моделей с учётом обратной связи от людей\,---\,парадигма, известная как RLHF.

% -----------------------------------------------------------
\section{Функция преимущества}
\label{sec:advantage}
% -----------------------------------------------------------

\subsection{Определение и основные свойства}

Вспомним из главы~\ref{ch:pg}, что теорема о градиенте стратегии выражает градиент целевой функции через функцию ценности действий:
\begin{equation}
\label{eq:pg-theorem-repeat}
\nabla_\theta J(\theta)
= \E_{S \sim d^\gamma_{\pi_\theta},\; A \sim \pi_\theta(\cdot|S)}
  \bigl[\nabla_\theta \log \pi_\theta(A \mid S)\; Q^{\pi_\theta}(S, A)\bigr].
\end{equation}
Глава~\ref{ch:pg} также показала, что любая \emph{зависящая от состояния} базовая линия $b(S_t)$ может быть вычтена из $Q^{\pi_\theta}(S_t, A_t)$ без изменения ожидаемого градиента (теорема~\ref{thm:baseline} в главе~\ref{ch:pg}). Наилучшей возможной базовой линией является та, которая максимально снижает дисперсию, и при мягких условиях оптимальным выбором является в точности функция ценности состояния $V^{\pi_\theta}(s)$. Остаток после вычитания этой базовой линии настолько важен, что заслуживает собственного имени.

\begin{definition}[Функция преимущества]
\label{def:advantage-ch9}
\index{advantage function}
Для стратегии $\pi$ и порождённых ею функций ценности действий и состояний $Q^\pi$ и $V^\pi$, \textbf{функция преимущества} определяется как
\begin{equation}
\label{eq:advantage-def}
A^\pi(s, a) \;=\; Q^\pi(s, a) - V^\pi(s).
\end{equation}
\end{definition}

Интуитивно, $A^\pi(s,a)$ измеряет относительное качество действия $a$ по сравнению со средним действием стратегии в состоянии $s$. Положительное преимущество означает, что $a$ лучше среднего; отрицательное \,---\, что хуже. Это свойство центрирования формализуется следующим утверждением.

\begin{proposition}[Нулевое среднее преимущества]
\label{prop:advantage-zero-mean}
\index{advantage function!zero-mean property}
Для любой стратегии $\pi$ и любого состояния $s \in \cS$,
\begin{equation}
\label{eq:advantage-zero-mean}
\E_{A \sim \pi(\cdot|s)}\bigl[A^\pi(s, A)\bigr] = 0.
\end{equation}
\end{proposition}

\begin{proof}
По определению~\ref{def:advantage-ch9} и линейности математического ожидания,
\[
\E_{A \sim \pi(\cdot|s)}\bigl[A^\pi(s, A)\bigr]
= \E_{A \sim \pi(\cdot|s)}\bigl[Q^\pi(s, A)\bigr] - V^\pi(s).
\]
Уравнение Беллмана (глава~\ref{ch:mdp}, уравнение~(3.12)) даёт
$V^\pi(s) = \sum_a \pi(a|s)\,Q^\pi(s,a)
           = \E_{A \sim \pi(\cdot|s)}[Q^\pi(s,A)]$,
поэтому правая часть равна нулю.
\end{proof}

Подстановка $V^{\pi_\theta}(S_t)$ вместо базовой линии в формулу градиента стратегии~\eqref{eq:pg-theorem-repeat} даёт \textbf{форму через преимущество} теоремы о градиенте стратегии:
\begin{equation}
\label{eq:advantage-pg}
\nabla_\theta J(\theta)
= \E_{S \sim d^\gamma_{\pi_\theta},\; A \sim \pi_\theta(\cdot|S)}
  \bigl[\nabla_\theta \log \pi_\theta(A \mid S)\; A^{\pi_\theta}(S, A)\bigr].
\end{equation}
Это версия теоремы о градиенте стратегии, лежащая в основе всех алгоритмов актор-критик.

\subsection{Почему преимущество снижает дисперсию}

Чтобы конкретно увидеть, почему использование $A^\pi(s,a)$ вместо $Q^\pi(s,a)$ помогает, рассмотрим простой пример.

\begin{example}[Преимущество против сырой отдачи]
\label{ex:advantage-variance}
Предположим, что агент находится в состоянии $s$, где каждое действие даёт высокую отдачу, поскольку $s$ уже является очень хорошим состоянием. Конкретно, пусть $Q^\pi(s, a) \in \{8, 10, 12\}$ с равными вероятностями для трёх действий, так что $V^\pi(s) = 10$. Оценки сырой функции ценности действий имеют среднее 10 и дисперсию $\Var[Q^\pi(s,\cdot)] = \frac{4+0+4}{3} \approx 2.67$.

Оценки преимущества равны $A^\pi(s,a) \in \{-2, 0, 2\}$, что имеет нулевое среднее (согласно утверждению~\ref{prop:advantage-zero-mean}) и ту же дисперсию 2.67. На первый взгляд дисперсия не изменилась, но обратите внимание на принципиальное различие: поскольку обновление градиента стратегии равно $\nabla_\theta \log \pi_\theta(a|s) \cdot A^\pi(s,a)$, обновления для посредственного действия ($A = 0$) равны нулю независимо от параметров стратегии. Использование сырого $Q$ давало бы ненулевой градиентный сигнал величиной 10 даже для нейтрального действия, засоряя обновление большим глобальным смещением, которое необходимо компенсировать на множестве выборок. Преимущество устраняет это смещение идеально и в ожидании.\qed
\end{example}

Более формально можно показать, что дисперсия оценщика градиента стратегии с оптимальной базовой линией $b^*(s)$ удовлетворяет
\begin{equation}
\label{eq:variance-reduction}
\Var\!\bigl[\nabla_\theta \log\pi_\theta(A|S)\,(Q^\pi(S,A) - b^*(S))\bigr]
\;\leq\;
\Var\!\bigl[\nabla_\theta \log\pi_\theta(A|S)\,Q^\pi(S,A)\bigr],
\end{equation}
причём снижение может быть значительным в средах, где функция ценности существенно варьируется по состояниям, но вариация преимущества внутри каждого состояния невелика.

% -----------------------------------------------------------
\section{Архитектура актор-критик}
\label{sec:actor-critic-arch}
% -----------------------------------------------------------

\subsection{Двухкомпонентный дизайн}

Алгоритм актор-критик поддерживает два взаимодействующих компонента.

\begin{definition}[Актор и критик]
\label{def:actor-critic}
\index{actor}\index{critic}
Алгоритм \textbf{актор-критик} поддерживает:
\begin{itemize}[leftmargin=*]
  \item \textbf{актора}: параметризованную стратегию $\pi_\theta(a \mid s)$, чьи параметры $\theta$ обновляются для максимизации $J(\theta)$;
  \item \textbf{критика}: параметризованную функцию ценности $V_\phi(s)$, чьи параметры $\phi$ обновляются для аппроксимации $V^{\pi_\theta}(s)$.
\end{itemize}
\end{definition}

Название отражает метафору исполнителя (актора) и рецензента (критика): после каждого действия критик сигнализирует, было ли это действие выше или ниже среднего, а актор корректируется соответственно.

На рисунке~\ref{fig:actor-critic-diagram} показан поток информации. Среда выдаёт состояние $S_t$; актор выбирает действие $A_t$; среда переходит в $S_{t+1}$ и выдаёт вознаграждение $R_{t+1}$. Критик использует $(S_t, R_{t+1}, S_{t+1})$ для вычисления ошибки TD $\delta_t$, которая служит одновременно и сигналом обучения критика (он измеряет, насколько ошибочным было предсказание самого критика), и оценкой преимущества для управления обновлением актора.


\begin{figure}[ht]
\centering
\begin{tikzpicture}[
    node distance=1.8cm and 2.8cm,
    box/.style={rectangle, draw, rounded corners=4pt, minimum width=2.4cm,
                minimum height=0.9cm, font=\small, fill=blue!8},
    env/.style={rectangle, draw, rounded corners=4pt, minimum width=2.4cm,
                minimum height=0.9cm, font=\small, fill=orange!12},
    arr/.style={-Stealth, thick},
    lab/.style={font=\footnotesize, inner sep=2pt}
  ]

  % Layout
  \node[box]  (actor)  at (0,0)    {Актор $\pi_\theta$};
  \node[env]  (env)    at (5.5,0)  {Среда};
  \node[box]  (critic) at (0,-2.5) {Критик $V_\phi$};

  % Actor -> Env: action A_t
  \draw[arr] ([yshift=-4pt]actor.east) -- ([yshift=-4pt]env.west)
        node[lab, below, midway] {действие $A_t$};

  % Env -> Actor: state S_t
  \draw[arr] ([yshift=4pt]env.west) -- ([yshift=4pt]actor.east)
        node[lab, above, midway] {состояние $S_t$};

  % Env -> Critic: reward + next state
  \draw[arr] (env.south) -- ++(0,-2.05) -- (critic.east)
        node[lab, above, pos=0.75] {$R_{t+1},\,S_{t+1}$};

  % Critic -> Actor: delta_t
  \draw[arr] (critic.north) -- (actor.south)
        node[lab, right, midway] {$\delta_t\!\approx\!\hat{A}_t$};

  % Annotation labels
  \node[lab, below=0.15cm of critic, text=blue!60!black] {\textit{обучает $V^{\pi_\theta}$}};
  \node[lab, above=0.15cm of actor, text=blue!60!black]  {\textit{обучает $\pi^*$}};

\end{tikzpicture}
\caption{Поток информации в алгоритме актор-критик. Актор выбирает действия; среда возвращает вознаграждения и следующие состояния; критик вырабатывает ошибку TD $\delta_t$, направляющую обновления как критика, так и актора.}
\label{fig:actor-critic-diagram}
\end{figure}

\subsection{Раздельные или совместные сети}

На практике актор и критик могут быть реализованы двумя способами. Если они являются \textbf{отдельными сетями}, актор $\pi_\theta$ и критик $V_\phi$ имеют полностью несвязанные параметры, что делает два целевых объекта обучения независимыми и часто проще поддающимися настройке. Если они \textbf{разделяют общий ствол} (общий набор скрытых слоёв), две головы разделяют представление, сокращая общее число параметров и стимулируя скрытые слои извлекать признаки, полезные как для стратегии, так и для оценки ценности. Общие стволы являются стандартом в современных реализациях, особенно для сложных наблюдаемых пространств\,---\,таких как изображения или токены естественного языка\,---\,где обучение хороших представлений само по себе является серьёзной задачей.

% -----------------------------------------------------------
\section{Ошибка TD как оценка преимущества}
\label{sec:td-error-advantage}
% -----------------------------------------------------------

\subsection{Одношаговая оценка}

Хотя форма через преимущество~\eqref{eq:advantage-pg} теоретически чиста, она требует знания как $Q^{\pi_\theta}$, так и $V^{\pi_\theta}$. Вычислять оба избыточно: поскольку $Q^\pi(s,a) = \E_\pi[R_{t+1} + \gamma V^\pi(S_{t+1}) \mid S_t=s, A_t=a]$, можно получить преимущество только через функцию ценности, используя наблюдаемый переход как выборку. Полученная величина является в точности одношаговой ошибкой TD, введённой в главе~\ref{ch:td}.

\begin{theorem}[Ошибка TD является несмещённой оценкой преимущества]
\label{thm:td-error-advantage}
\index{TD error!as advantage estimate}
Пусть $V_\phi = V^{\pi_\theta}$ (критик идеален). Тогда для любого перехода $(S_t, A_t, R_{t+1}, S_{t+1})$, выборочного согласно $\pi_\theta$,
\begin{equation}
\label{eq:td-as-advantage}
\E\bigl[\delta_t \;\big|\; S_t = s,\; A_t = a\bigr]
\;=\; A^{\pi_\theta}(s, a),
\end{equation}
где $\delta_t = R_{t+1} + \gamma V^{\pi_\theta}(S_{t+1}) - V^{\pi_\theta}(S_t)$.
\end{theorem}

\begin{proof}
Взяв условное математическое ожидание $\delta_t$ при $S_t = s$ и $A_t = a$:
\begin{align*}
\E[\delta_t \mid S_t = s, A_t = a]
&= \E\bigl[R_{t+1} + \gamma V^{\pi_\theta}(S_{t+1}) \mid S_t=s, A_t=a\bigr]
   - V^{\pi_\theta}(s).
\end{align*}
По уравнению Беллмана для $Q^{\pi_\theta}$ (глава~\ref{ch:mdp}),
\[
\E\bigl[R_{t+1} + \gamma V^{\pi_\theta}(S_{t+1}) \mid S_t=s, A_t=a\bigr]
= Q^{\pi_\theta}(s, a),
\]
что следует из того, что марковское свойство позволяет разложить $\E[G_{t+1} \mid S_t, A_t] = \E[\E[G_{t+1} \mid S_{t+1}] \mid S_t, A_t] = \E[V^{\pi_\theta}(S_{t+1}) \mid S_t, A_t]$.
Поэтому
\[
\E[\delta_t \mid S_t=s, A_t=a]
= Q^{\pi_\theta}(s,a) - V^{\pi_\theta}(s)
= A^{\pi_\theta}(s,a). \qedhere
\]
\end{proof}

Теорема~\ref{thm:td-error-advantage} является ключевым результатом, делающим алгоритмы актор-критик практичными. Вместо поддержки отдельных аппроксиматоров функций для $Q^\pi$ и $V^\pi$, достаточно лишь обучить $V_\phi$; ошибка TD $\delta_t$ тогда является одновыборочной несмещённой оценкой преимущества.

\begin{remark}
Несмещённость выполняется только при $V_\phi = V^{\pi_\theta}$. На практике критик несовершенен, поэтому ошибка TD вносит \emph{смещение} в оценку преимущества. Компромисс между смещением и дисперсией является центральным противоречием в методах актор-критик: хорошо обученный критик существенно снижает дисперсию, но вносит систематические ошибки. Это то же противоречие, которое отличает TD(0) от Монте-Карло в главе~\ref{ch:td}.
\end{remark}

\subsection{Одношаговый алгоритм актор-критик}

Подстановка $\delta_t$ в градиент стратегии в форме преимущества даёт простейшее обновление актор-критик с обучением на политике.

\begin{algorithm}[ht]
\caption{Одношаговый актор-критик с обучением на политике (критик TD(0))}
\label{alg:ac-td0}
\DontPrintSemicolon
\KwIn{Сеть стратегии $\pi_\theta$, сеть ценности $V_\phi$, размеры шагов $\alpha_\theta, \alpha_\phi > 0$, дисконт $\gamma$}
Инициализировать $\theta$ и $\phi$ случайно\;
\For{каждый эпизод}{
  Наблюдать начальное состояние $S_0$\;
  \For{$t = 0, 1, 2, \ldots$ до терминального состояния}{
    Выбрать действие $A_t \sim \pi_\theta(\cdot \mid S_t)$\;
    Наблюдать $R_{t+1}$ и $S_{t+1}$\;
    Вычислить ошибку TD:
    $\delta_t \leftarrow R_{t+1} + \gamma V_\phi(S_{t+1}) - V_\phi(S_t)$\;
    \tcp{Обновление критика: полуградиент TD(0)}
    $\phi \;\leftarrow\; \phi + \alpha_\phi\,\delta_t\,\nabla_\phi V_\phi(S_t)$\;
    \tcp{Обновление актора: градиент стратегии с оценкой преимущества через ошибку TD}
    $\theta \;\leftarrow\; \theta
      + \alpha_\theta\,\gamma^t\,\delta_t\,
        \nabla_\theta \log \pi_\theta(A_t \mid S_t)$\;
  }
}
\end{algorithm}

Обновление критика в алгоритме~\ref{alg:ac-td0} является в точности полуградиентным обновлением TD(0), выведенным в главе~\ref{ch:td} (уравнение~(6.18)). Обновление актора заменяет полную отдачу $G_t$ REINFORCE на ошибку TD $\delta_t$, которая доступна сразу после каждого шага и имеет значительно меньшую дисперсию. Принципиально важно, что ни одно из обновлений не требует ожидания конца эпизода, что делает методы актор-критик применимыми к непрерывным (не эпизодическим) задачам.

\begin{example}[CartPole с одношаговым актор-критиком]
\label{ex:cartpole-ac}
Рассмотрим среду CartPole, где состоянием является четырёхмерный вектор (положение тележки, скорость тележки, угол шеста, угловая скорость шеста), действие является бинарным (толкнуть влево или вправо), а эпизод завершается, когда шест падает или тележка выезжает за пределы трассы.

Актор \,---\, это двухслойная сеть с выходом softmax для двух действий; критик \,---\, двухслойная сеть, выдающая скаляр $V_\phi(s)$. После шага, на котором шест почти упал, а затем восстановился,
\[
R_{t+1} = 1,\quad V_\phi(S_t) = 4.2,\quad V_\phi(S_{t+1}) = 3.1,
\]
поэтому $\delta_t = 1 + 0.99 \times 3.1 - 4.2 = -0.131$. Отрицательная ошибка TD сигнализирует, что состояние $S_t$ было переоценено\,---\,агент действует хуже, чем ожидалось,\,---\,и обновление актора слегка снижает вероятность действия, совершённого в $S_t$. Напротив, на шаге, где шест стабилен и критик недооценил ситуацию, $\delta_t > 0$ повышает вероятность стабилизирующего действия. Эта немедленная, пошаговая обратная связь отличает актор-критик от REINFORCE.\qed
\end{example}

% -----------------------------------------------------------
\section{Обобщённая оценка преимущества}
\label{sec:gae-ch9}
% -----------------------------------------------------------

\subsection{Многошаговые отдачи и $n$-шаговое преимущество}

Одношаговая оценка преимущества $\delta_t$ имеет низкую дисперсию, но смещена, когда $V_\phi \neq V^{\pi_\theta}$. На другом крайнем случае, оценка преимущества Монте-Карло $G_t - V_\phi(S_t)$ является несмещённой оценкой преимущества, но имеет высокую дисперсию, поскольку суммирует все будущие вознаграждения. Глава~\ref{ch:td} показала, что $n$-шаговые отдачи интерполируют между этими полюсами: $n$-шаговая цель TD равна
\[
G_t^{(n)} = R_{t+1} + \gamma R_{t+2} + \cdots + \gamma^{n-1}R_{t+n}
             + \gamma^n V_\phi(S_{t+n}).
\]
Соответствующая $n$-шаговая оценка преимущества:
\begin{equation}
\label{eq:n-step-advantage}
\hat{A}_t^{(n)}
= G_t^{(n)} - V_\phi(S_t)
= \sum_{l=0}^{n-1} \gamma^l R_{t+l+1} + \gamma^n V_\phi(S_{t+n})
  - V_\phi(S_t).
\end{equation}
При $n \to \infty$ $\hat{A}_t^{(n)} \to G_t - V_\phi(S_t)$\,---\,оценке преимущества Монте-Карло. При $n = 1$ получаем $\delta_t$.

Вместо фиксации конкретного $n$, Обобщённая Оценка Преимущества (GAE), введённая Шульманом и~др.~(2016), берёт экспоненциально взвешенное среднее всех $n$-шаговых оценок.

\subsection{Определение и вывод GAE}

\begin{definition}[Обобщённая оценка преимущества]
\label{def:gae-ch9}
\index{Generalized Advantage Estimation (GAE)}
\index{GAE|see{Generalized Advantage Estimation}}
Для параметров $\gamma \in (0,1)$ и $\lambda \in [0,1]$ определим ошибки TD
\[
\delta_t = R_{t+1} + \gamma V_\phi(S_{t+1}) - V_\phi(S_t),
\]
и \textbf{оценку преимущества GAE} как
\begin{equation}
\label{eq:gae-def}
\hat{A}_t^{\mathrm{GAE}(\gamma,\lambda)}
= \sum_{l=0}^{T-t-1} (\gamma\lambda)^l\,\delta_{t+l}.
\end{equation}
\end{definition}

Формула~\eqref{eq:gae-def} выглядит как экспоненциально затухающая сумма одношаговых ошибок TD, начиная с момента $t$. Она эквивалентна $\mathrm{TD}(\lambda)$ для оценки ценности, применённому здесь для оценки преимущества. Следующее утверждение явно устанавливает связь с $n$-шаговыми преимуществами.

\begin{proposition}[GAE как взвешенное среднее $n$-шаговых преимуществ]
\label{prop:gae-nstep}
Оценка GAE равна нормализованной на $(1-\lambda)$ смеси всех $n$-шаговых оценок преимущества:
\begin{equation}
\label{eq:gae-nstep}
\hat{A}_t^{\mathrm{GAE}(\gamma,\lambda)}
= (1-\lambda)\sum_{n=1}^{T-t} \lambda^{n-1}\,\hat{A}_t^{(n)}.
\end{equation}
\end{proposition}

\begin{proof}
Раскроем $\hat{A}_t^{(n)}$ из~\eqref{eq:n-step-advantage} через ошибки TD. Заметим, что
\[
\hat{A}_t^{(n)}
= \sum_{l=0}^{n-1}\gamma^l R_{t+l+1} + \gamma^n V_\phi(S_{t+n}) - V_\phi(S_t).
\]
Добавляя и вычитая $\gamma^l V_\phi(S_{t+l})$ для каждого промежуточного шага,
\begin{align*}
\hat{A}_t^{(n)}
&= \sum_{l=0}^{n-1} \gamma^l
   \bigl(R_{t+l+1} + \gamma V_\phi(S_{t+l+1}) - V_\phi(S_{t+l})\bigr)
 = \sum_{l=0}^{n-1}\gamma^l\,\delta_{t+l}.
\end{align*}
Теперь формируем взвешенную сумму:
\begin{align*}
(1-\lambda)\sum_{n=1}^{\infty}\lambda^{n-1}\hat{A}_t^{(n)}
&= (1-\lambda)\sum_{n=1}^{\infty}\lambda^{n-1}
   \sum_{l=0}^{n-1}\gamma^l\,\delta_{t+l} \\
&= (1-\lambda)\sum_{l=0}^{\infty}(\gamma\lambda)^l\,\delta_{t+l}
   \underbrace{\sum_{n=l+1}^{\infty}(1-\lambda)\lambda^{n-1-l}}_{=\,1} \\
&= \sum_{l=0}^{\infty}(\gamma\lambda)^l\,\delta_{t+l},
\end{align*}
где перестановка порядка суммирования допустима для конечно-горизонтных эпизодов (внутренняя сумма конечна), а заключённый в скобки геометрический ряд суммируется до 1, поскольку $(1-\lambda)\sum_{k=0}^{\infty}\lambda^k = 1$. Усечение до $l = T-t-1$ даёт~\eqref{eq:gae-nstep}.
\end{proof}

\subsection{Параметр $\lambda$: компромисс между смещением и дисперсией}

Роль $\lambda$ состоит в интерполяции между двумя крайними случаями.

\begin{itemize}[leftmargin=*]
  \item \textbf{$\lambda = 0$}: В~\eqref{eq:gae-def} выживает только первый член, поэтому $\hat{A}_t^{\mathrm{GAE}(\gamma,0)} = \delta_t$. Это одношаговая ошибка TD\,---\,минимальная дисперсия, но смещённая, когда $V_\phi \neq V^{\pi_\theta}$.
  \item \textbf{$\lambda = 1$}: Сумма простирается до конца эпизода и $\hat{A}_t^{\mathrm{GAE}(\gamma,1)} = \sum_{l=0}^{T-t-1}\gamma^l \delta_{t+l} = G_t - V_\phi(S_t)$\,---\,оценка преимущества Монте-Карло. Это несмещённо (для любого $V_\phi$), но с максимальной дисперсией.
\end{itemize}

При промежуточных $\lambda$ эффективный горизонт убывает экспоненциально: члены $\delta_{t+l}$ взвешиваются с весом $(\gamma\lambda)^l$, поэтому значимо вносят вклад лишь примерно первые $1/(1-\gamma\lambda)$ шагов. Значения $\lambda \in [0.90, 0.97]$ при $\gamma \approx 0.99$ типичны в современных реализациях, обеспечивая хороший баланс для эпизодов умеренной длины.

На рисунке~\ref{fig:gae-lambda} визуализируется, как изменяется эффективный вес при ошибке TD на шаге $l$ в зависимости от $\lambda$.

\begin{figure}[ht]
\centering
\begin{tikzpicture}[xscale=0.55, yscale=3.5]
  % Оси
  \draw[->] (0,0) -- (12,0) node[right, font=\small] {смещение шага $l$};
  \draw[->] (0,0) -- (0,1.15) node[above, font=\small] {вес $(\gamma\lambda)^l$};

  % lambda = 0: всплеск только при l=0
  \fill[blue!60] (0,0) rectangle (0.35,1.0);

  % lambda = 0.5 bars
  \fill[orange!70] (0.37,0) rectangle (0.72,0.990);
  \fill[orange!70] (1.37,0) rectangle (1.72,0.495);
  \fill[orange!70] (2.37,0) rectangle (2.72,0.248);
  \fill[orange!70] (3.37,0) rectangle (3.72,0.124);
  \fill[orange!70] (4.37,0) rectangle (4.72,0.062);
  \fill[orange!70] (5.37,0) rectangle (5.72,0.031);
  \fill[orange!70] (6.37,0) rectangle (6.72,0.015);

  % lambda = 0.95 bars
  \fill[green!60!black,opacity=0.8] (0.74,0) rectangle (1.09,0.990);
  \fill[green!60!black,opacity=0.8] (1.74,0) rectangle (2.09,0.941);
  \fill[green!60!black,opacity=0.8] (2.74,0) rectangle (3.09,0.894);
  \fill[green!60!black,opacity=0.8] (3.74,0) rectangle (4.09,0.849);
  \fill[green!60!black,opacity=0.8] (4.74,0) rectangle (5.09,0.807);
  \fill[green!60!black,opacity=0.8] (5.74,0) rectangle (6.09,0.767);
  \fill[green!60!black,opacity=0.8] (6.74,0) rectangle (7.09,0.728);
  \fill[green!60!black,opacity=0.8] (7.74,0) rectangle (8.09,0.692);
  \fill[green!60!black,opacity=0.8] (8.74,0) rectangle (9.09,0.657);
  \fill[green!60!black,opacity=0.8] (9.74,0) rectangle (10.09,0.624);

  % lambda = 1 bars (равномерный вес gamma^l ~ 0.99^l ~ константа для короткого эп.)
  \fill[red!60,opacity=0.7] (1.11,0) rectangle (1.46,0.990);
  \fill[red!60,opacity=0.7] (2.11,0) rectangle (2.46,0.990);
  \fill[red!60,opacity=0.7] (3.11,0) rectangle (3.46,0.990);
  \fill[red!60,opacity=0.7] (4.11,0) rectangle (4.46,0.990);
  \fill[red!60,opacity=0.7] (5.11,0) rectangle (5.46,0.990);
  \fill[red!60,opacity=0.7] (6.11,0) rectangle (6.46,0.990);
  \fill[red!60,opacity=0.7] (7.11,0) rectangle (7.46,0.990);
  \fill[red!60,opacity=0.7] (8.11,0) rectangle (8.46,0.990);
  \fill[red!60,opacity=0.7] (9.11,0) rectangle (9.46,0.990);
  \fill[red!60,opacity=0.7] (10.11,0) rectangle (10.46,0.990);

  % Отметки по оси y
  \draw (0.05,1) -- (-0.05,1) node[left, font=\tiny] {1};
  \draw (0.05,0.5) -- (-0.05,0.5) node[left, font=\tiny] {0.5};

  % Легенда внутри
  \node[fill=blue!60, minimum width=0.25cm, minimum height=0.2cm,
        inner sep=0pt] at (2.0, 1.05) {};
  \node[right, font=\tiny] at (2.15, 1.05) {$\lambda=0$};
  \node[fill=orange!70, minimum width=0.25cm, minimum height=0.2cm,
        inner sep=0pt] at (4.0, 1.05) {};
  \node[right, font=\tiny] at (4.15, 1.05) {$\lambda=0.5$};
  \node[fill=green!60!black, minimum width=0.25cm, minimum height=0.2cm,
        inner sep=0pt, opacity=0.8] at (6.0, 1.05) {};
  \node[right, font=\tiny] at (6.15, 1.05) {$\lambda=0.95$};
  \node[fill=red!60, minimum width=0.25cm, minimum height=0.2cm,
        inner sep=0pt, opacity=0.7] at (8.2, 1.05) {};
  \node[right, font=\tiny] at (8.35, 1.05) {$\lambda=1$};

\end{tikzpicture}
\caption{Эффективные веса $(\gamma\lambda)^l$ при ошибках TD с шаговым смещением $l$ в GAE, для $\gamma = 0.99$ и четырёх значений $\lambda$. При $\lambda=0$ только непосредственная ошибка TD вносит вклад; при $\lambda=1$ все шаги получают почти равномерный вес (убывающий лишь с коэффициентом $\gamma^l \approx 1$). Промежуточные $\lambda$ дают мягкое экспоненциальное убывание.}
\label{fig:gae-lambda}
\end{figure}

\subsection{Цель GAE для критика}

GAE также предоставляет естественную цель для обучения критика. Поскольку $\hat{A}_t^{\mathrm{GAE}} = \hat{V}_t - V_\phi(S_t)$, где
\begin{equation}
\label{eq:gae-vtarget}
\hat{V}_t = \hat{A}_t^{\mathrm{GAE}} + V_\phi(S_t),
\end{equation}
критик может быть обучен путём минимизации среднеквадратичной ошибки $\E[(V_\phi(S_t) - \hat{V}_t)^2]$ при фиксированной цели $\hat{V}_t$. Это создаёт согласованную пару: актор использует $\hat{A}_t^{\mathrm{GAE}}$ для градиента стратегии, а критик использует $\hat{V}_t$ как цель регрессии.

% -----------------------------------------------------------
\section{A2C и A3C}
\label{sec:a3c}
% -----------------------------------------------------------

\subsection{Параллелизм как инструмент снижения дисперсии}

Одной из практических трудностей алгоритмов актор-критик с обучением на политике является то, что последовательные переходы $\{(S_t, A_t, R_{t+1}, S_{t+1})\}$, собранные одним агентом, сильно \emph{коррелированы}: агент следует единственной цепи Маркова, поэтому соседние выборки имеют схожие состояния. Высокая корреляция увеличивает эффективную дисперсию оценки градиента даже при низкой дисперсии отдельных выборок.

Мних и~др.~(2016) предложили \textbf{Асинхронный Актор-Критик с Преимуществом} (A3C)\index{A3C} в качестве решения. Идея состоит в запуске нескольких независимых \emph{рабочих} процессов, каждый со своей копией среды и локальной копией параметров сети. Рабочие генерируют опыт независимо и асинхронно отправляют обновления градиента центральному серверу параметров, который агрегирует обновления.

\begin{definition}[Архитектура A3C]
\label{def:a3c}
\index{A3C}
В A3C $K$ рабочих процессов выполняются параллельно. Каждый рабочий $k$:
\begin{enumerate}[label=\arabic*., leftmargin=*]
  \item Копирует глобальные параметры $(\theta, \phi)$ в локальные копии $(\theta_k, \phi_k)$.
  \item Собирает короткую траекторию из $t_\mathrm{max}$ переходов с использованием $\pi_{\theta_k}$.
  \item Вычисляет градиенты актора и критика с помощью оценки преимущества $\hat{A}_t$ из~\eqref{eq:n-step-advantage}.
  \item Асинхронно применяет градиенты к глобальным $(\theta, \phi)$.
\end{enumerate}
\end{definition}

Поскольку разные рабочие, как правило, оказываются в разных состояниях и исследуют разные части пространства состояний одновременно, их градиенты почти некоррелированы даже без буфера воспроизведения опыта. Эта декорреляция является ключевым преимуществом асинхронного параллелизма.

\subsection{A2C: синхронный вариант}

\textbf{Актор-Критик с Преимуществом} (A2C)\index{A2C} \,---\, это синхронная версия A3C. Все $K$ рабочих одновременно собирают пакет переходов, и обновление параметров выполняется только после того, как все рабочие отчитались. Поскольку обновления синхронны, нет несоответствия между параметрами, используемыми для генерации данных, и теми, которые обновляются.

\begin{algorithm}[ht]
\caption{Актор-Критик с Преимуществом (A2C)}
\label{alg:a2c}
\DontPrintSemicolon
\KwIn{$K$ параллельных рабочих, длина прокрутки $T_\mathrm{roll}$, дисконт $\gamma$, параметр GAE $\lambda$}
Инициализировать глобальные $\theta, \phi$\;
\Repeat{сходимость}{
  \For{каждый рабочий $k = 1, \ldots, K$ \emph{параллельно}}{
    Запустить $\pi_\theta$ на $T_\mathrm{roll}$ шагах; собрать $\{(S_t^k, A_t^k, R_{t+1}^k)\}_{t=0}^{T_\mathrm{roll}-1}$\;
    Вычислить $\hat{A}_t^k$ через GAE~\eqref{eq:gae-def}\;
    Вычислить $\hat{V}_t^k$ через~\eqref{eq:gae-vtarget}\;
  }
  Агрегировать градиенты по всем рабочим\;
  Обновить $\theta$ с помощью градиента актора $\frac{1}{KT_\mathrm{roll}}\sum_k\sum_t \hat{A}_t^k\,\nabla_\theta \log\pi_\theta(A_t^k|S_t^k)$\;
  Обновить $\phi$ минимизацией $\frac{1}{KT_\mathrm{roll}}\sum_k\sum_t (V_\phi(S_t^k) - \hat{V}_t^k)^2$\;
}
\end{algorithm}

На практике A2C достигает производительности, сопоставимой с A3C, и проще воспроизводимо реализуется, поскольку синхронный дизайн устраняет состояния гонки. Современные пайплайны глубокого ОП почти повсеместно используют синхронное пакетирование в стиле A2C в качестве основы, дополняемой ограничениями доверительной области, разработанными в следующем разделе.

% -----------------------------------------------------------
\section{Методы доверительной области}
\label{sec:trpo}
% -----------------------------------------------------------

\subsection{Повторное использование данных с другой политики через взвешивание по важности}

Методы с обучением на политике, такие как A2C, собирают пакет переходов и затем немедленно отбрасывают их после одного шага градиента. Это расточительно: каждое взаимодействие со средой обходится дорого (по времени симуляции, вычислениям или реальным затратам), но ограничение обучения на политике вынуждает отбрасывать каждую выборку после изменения параметров.

Естественным средством является взвешивание по важности (введённое в главе~\ref{ch:mc} в контексте Монте-Карло с данными другой политики). Предположим, что данные были собраны с использованием \emph{старой} стратегии $\pi_{\theta_\mathrm{old}}$, и мы хотим оценить или оптимизировать \emph{новую} стратегию $\pi_\theta$. Определим пошаговое \textbf{отношение важности}:\index{importance ratio}
\begin{equation}
\label{eq:importance-ratio}
r_t(\theta)
= \frac{\pi_\theta(A_t \mid S_t)}{\pi_{\theta_\mathrm{old}}(A_t \mid S_t)}.
\end{equation}
По стандартному взвешиванию по важности, градиент стратегии $\E_{\pi_\theta}[\nabla_\theta \log\pi_\theta(A|S)\,A^\pi(S,A)]$ можно переписать, при приближении, что распределение состояний, индуцированное $\pi_\theta$, близко к распределению $\pi_{\theta_\mathrm{old}}$, как
\begin{equation}
\label{eq:surrogate-obj}
L(\theta) = \E_{\pi_{\theta_\mathrm{old}}}\bigl[r_t(\theta)\,\hat{A}_t\bigr].
\end{equation}
Функция $L(\theta)$ называется \textbf{суррогатной целевой функцией}.
\index{surrogate objective}
Её можно оптимизировать с использованием данных из $\pi_{\theta_\mathrm{old}}$, что позволяет выполнять несколько шагов градиента на одном пакете.

\begin{proposition}[Суррогатный градиент равен градиенту стратегии при $\theta_\mathrm{old}$]
\label{prop:surrogate-gradient}
\begin{equation}
\label{eq:surrogate-gradient}
\nabla_\theta L(\theta)\big|_{\theta = \theta_\mathrm{old}}
= \nabla_\theta J(\theta)\big|_{\theta = \theta_\mathrm{old}}.
\end{equation}
\end{proposition}

\begin{proof}
При $\theta = \theta_\mathrm{old}$, $r_t(\theta) = 1$ для всех переходов. Дифференцируя $L(\theta) = \E[r_t(\theta)\,\hat{A}_t]$, получаем
\[
\nabla_\theta L(\theta)\big|_{\theta_\mathrm{old}}
= \E\bigl[\nabla_\theta r_t(\theta)\big|_{\theta_\mathrm{old}}\cdot \hat{A}_t\bigr].
\]
Поскольку $\nabla_\theta r_t(\theta) = \frac{\nabla_\theta \pi_\theta(A_t|S_t)}{\pi_{\theta_\mathrm{old}}(A_t|S_t)}$, при $\theta = \theta_\mathrm{old}$ это равно $\nabla_\theta \log \pi_\theta(A_t|S_t)|_{\theta_\mathrm{old}}$, что даёт градиент стратегии.
\end{proof}

\subsection{Опасность больших обновлений}

Утверждение~\ref{prop:surrogate-gradient} показывает, что суррогатная целевая функция локально корректна, но что происходит, когда $\theta$ значительно отклоняется от $\theta_\mathrm{old}$?

\begin{proposition}[Граница монотонного улучшения]
\label{prop:monotone-improvement}
\index{trust region!monotone improvement bound}
Существует константа $C > 0$, зависящая от $\gamma$ и $\max_{s,a}|A^{\pi_\theta}(s,a)|$, такая что
\begin{equation}
\label{eq:improvement-bound}
J(\theta) \geq L(\theta)
- C \cdot \max_s \KL{\pi_\theta(\cdot|s)}{\pi_{\theta_\mathrm{old}}(\cdot|s)}.
\end{equation}
\end{proposition}

Граница~\eqref{eq:improvement-bound}, принадлежащая Шульману и~др.~(2015), говорит, что $J(\theta)$ ограничена снизу $L(\theta)$ минус член, пропорциональный KL-дивергенции новой стратегии от старой. Если KL мала, суррогат $L(\theta)$ является надёжной аппроксимацией $J(\theta)$, поэтому максимизация $L$ гарантирует улучшение $J$. Если KL велика, граница неинформативна: суррогат может расти, пока истинная целевая функция падает.

Это фундаментальная мотивация подхода доверительной области\index{trust region}: мы должны двигать $\theta$ только в области, где KL-дивергенция от $\theta_\mathrm{old}$ достаточно мала, чтобы суррогат был надёжным.

\subsection{Trust Region Policy Optimization (TRPO)}

\textbf{Trust Region Policy Optimization} (TRPO)\index{TRPO}, введённый Шульманом и~др.~(2015), решает задачу оптимизации с ограничениями:
\begin{equation}
\label{eq:trpo}
\max_\theta\; L(\theta)
\quad\mbox{при ограничении}\quad
\max_s\;\KL{\pi_\theta(\cdot|s)}{\pi_{\theta_\mathrm{old}}(\cdot|s)} \leq \delta,
\end{equation}
где $\delta > 0$ \,---\, гиперпараметр, управляющий размером доверительной области. На практике максимум по состояниям заменяется эмпирическим средним.

Решение~\eqref{eq:trpo} требует вычисления матрицы информации Фишера (гессиана ограничения KL) и решения задачи квадратичного программирования с ограничениями. Это приближённо выполняется методом сопряжённых градиентов с поиском по направлению, однако процедура вычислительно тяжела: она требует нескольких произведений гессиана на вектор за обновление, что плохо масштабируется для сетей с миллионами параметров. TRPO теоретически хорошо обоснован, но редко используется на практике для больших моделей.

PPO, разработанный далее, предоставляет значительно более простое приближение к той же идее.

% -----------------------------------------------------------
\section{Proximal Policy Optimization}
\label{sec:ppo}
% -----------------------------------------------------------

\subsection{От доверительных областей к отсечению}

Proximal Policy Optimization (PPO)\index{PPO}, введённый Шульманом и~др.~(2017), заменяет явное ограничение KL TRPO простой операцией отсечения отношения важности. Интуиция состоит в том, что $r_t(\theta) = 1$ при $\theta = \theta_\mathrm{old}$, и что отклонение $r_t(\theta)$ далеко от 1 соответствует большому изменению стратегии. PPO непосредственно предотвращает большие отклонения, ограничивая $r_t$ интервалом $[1-\varepsilon, 1+\varepsilon]$ для гиперпараметра $\varepsilon > 0$ (обычно $\varepsilon \in [0.1, 0.3]$).

\begin{definition}[Отсечённая целевая функция PPO]
\label{def:ppo-clip}
\index{PPO!clipped objective}
\textbf{Отсечённая суррогатная целевая функция PPO} определяется как
\begin{equation}
\label{eq:ppo-clip}
L^{\mathrm{CLIP}}(\theta)
= \E\!\left[\min\!\Bigl(
    r_t(\theta)\,\hat{A}_t,\;
    \operatorname{clip}\!\bigl(r_t(\theta),\,1-\varepsilon,\,1+\varepsilon\bigr)
    \,\hat{A}_t
  \Bigr)\right],
\end{equation}
где $\operatorname{clip}(x, a, b) = \max(a, \min(b, x))$.
\end{definition}

\subsection{Анализ случаев механизма отсечения}

\begin{theorem}[Корректность отсечения PPO]
\label{thm:ppo-clip-cases}
\index{PPO!clipping cases}
Для каждого перехода отсечённая целевая функция $L^{\mathrm{CLIP}}$ обладает следующими свойствами:
\begin{enumerate}[label=(\roman*), leftmargin=*]
  \item Если $\hat{A}_t > 0$: $L^{\mathrm{CLIP}}$ неубывает по $r_t(\theta)$ при $r_t \leq 1+\varepsilon$ и постоянна при $r_t > 1+\varepsilon$. Градиент равен нулю при $r_t > 1+\varepsilon$.
  \item Если $\hat{A}_t < 0$: $L^{\mathrm{CLIP}}$ невозрастает по $r_t(\theta)$ при $r_t \geq 1-\varepsilon$ и постоянна при $r_t < 1-\varepsilon$. Градиент равен нулю при $r_t < 1-\varepsilon$.
  \item При $\theta = \theta_\mathrm{old}$ (т.\,е. $r_t = 1 \in [1-\varepsilon, 1+\varepsilon]$) градиент $L^{\mathrm{CLIP}}$ равен градиенту стратегии.
\end{enumerate}
\end{theorem}

\begin{proof}
Зафиксируем переход и запишем $r = r_t(\theta)$, $A = \hat{A}_t$. Пусть $f_1(r) = r A$ и $f_2(r) = \operatorname{clip}(r, 1-\varepsilon, 1+\varepsilon)\cdot A$.

\medskip
\textit{Случай (i): $A > 0$.}
Тогда $f_1$ строго возрастает по $r$. При $r \leq 1+\varepsilon$ отсечение неактивно: $f_2(r) = rA = f_1(r)$. При $r > 1+\varepsilon$, $f_2(r) = (1+\varepsilon)A$, константа. Таким образом,
\[
\min(f_1, f_2)
= \begin{cases} rA & r \leq 1+\varepsilon, \\ (1+\varepsilon)A & r > 1+\varepsilon.\end{cases}
\]
Это неубывающая функция, и её производная по $r$ равна нулю при $r > 1+\varepsilon$. Не существует стимула в виде градиента для увеличения вероятности хорошего действия сверх предела $(1+\varepsilon)$.

\medskip
\textit{Случай (ii): $A < 0$.}
Тогда $f_1$ строго убывает. При $r \geq 1-\varepsilon$, $f_2(r) = rA = f_1(r)$. При $r < 1-\varepsilon$, $f_2(r) = (1-\varepsilon)A$, константа, и $f_1(r) > f_2(r)$, поскольку $r < 1-\varepsilon$ и $A<0$.
Поэтому $\min(f_1,f_2) = f_2 = (1-\varepsilon)A$ при $r < 1-\varepsilon$ (постоянна), и $\min(f_1,f_2) = rA$ при $r \geq 1-\varepsilon$. Следовательно:
\[
\min(f_1, f_2)
= \begin{cases} (1-\varepsilon)A & r < 1-\varepsilon, \\ rA & r \geq 1-\varepsilon.\end{cases}
\]
Градиент по $r$ равен нулю при $r < 1-\varepsilon$: нет стимула агрессивно снижать вероятность плохого действия ниже пола $(1-\varepsilon)$.

\medskip
\textit{Случай (iii):} При $r = 1 \in (1-\varepsilon, 1+\varepsilon)$ отсечение неактивно в обоих случаях, поэтому $L^{\mathrm{CLIP}} = \E[r_t(\theta)\hat{A}_t] = L(\theta)$ в окрестности $\theta_\mathrm{old}$. По утверждению~\ref{prop:surrogate-gradient}, $\nabla_\theta L^{\mathrm{CLIP}} |_{\theta_\mathrm{old}} = \nabla_\theta J(\theta)|_{\theta_\mathrm{old}}$.
\end{proof}

Оператор $\min$ в определении~\ref{def:ppo-clip} является принципиальным: он делает $L^{\mathrm{CLIP}}$ \emph{пессимистической нижней границей} суррогата $r_t(\theta)\hat{A}_t$. Рисунок~\ref{fig:ppo-clip} иллюстрирует два случая.

\begin{figure}[ht]
\centering
\begin{tikzpicture}[xscale=3.5, yscale=2.2]
  % Случай 1: A > 0
  \begin{scope}[xshift=0cm]
    \def\eps{0.2}
    % Оси
    \draw[->] (0.5,0) -- (2.0,0) node[right, font=\small] {$r_t(\theta)$};
    \draw[->] (0.6,-0.1) -- (0.6,1.1)
              node[above, font=\small] {$L^{\mathrm{CLIP}}$};

    % Необрезанный суррогат r*A (пунктир)
    \draw[dashed, gray, thick] (0.65, 0.325) -- (1.9, 0.95);
    \node[gray, font=\tiny, right] at (1.85, 0.90) {$r\hat{A}$};

    % Отсечённая целевая функция
    \draw[blue, very thick] (0.65, 0.325) -- (1.2, 0.6);
    \draw[blue, very thick] (1.2, 0.6) -- (1.9, 0.6);
    \fill[blue] (1.2, 0.6) circle (1.2pt);

    % Отметки 1-eps, 1, 1+eps на оси x
    \draw (0.8, 0.02) -- (0.8,-0.02) node[below, font=\tiny] {$1{-}\varepsilon$};
    \draw (1.0, 0.02) -- (1.0,-0.02) node[below, font=\tiny] {$1$};
    \draw (1.2, 0.02) -- (1.2,-0.02) node[below, font=\tiny] {$1{+}\varepsilon$};

    % Вертикальная пунктирная линия при 1+eps
    \draw[dashed, gray!60] (1.2, 0) -- (1.2, 0.6);

    \node[font=\small, align=center] at (1.25, -0.28)
          {$\hat{A}_t > 0$: отсечение при $1+\varepsilon$};
  \end{scope}

  % Случай 2: A < 0
  \begin{scope}[xshift=2.3cm]
    % Оси
    \draw[->] (0.5,0) -- (2.0,0) node[right, font=\small] {$r_t(\theta)$};
    \draw[->] (0.6,-0.85) -- (0.6,0.15)
              node[above, font=\small] {$L^{\mathrm{CLIP}}$};

    % Необрезанный суррогат r*A (пунктир, отрицательный наклон)
    \draw[dashed, gray, thick] (0.65, -0.065) -- (1.9, -0.57);
    \node[gray, font=\tiny, right] at (1.85, -0.55) {$r\hat{A}$};

    % Отсечённая целевая функция
    \draw[red!80!black, very thick] (0.65, -0.325) -- (0.8, -0.325);
    \draw[red!80!black, very thick] (0.8, -0.325) -- (1.9, -0.57);
    \fill[red!80!black] (0.8, -0.325) circle (1.2pt);

    % Отметки 1-eps, 1, 1+eps на оси x
    \draw (0.8, 0.015) -- (0.8,-0.015) node[above, font=\tiny] {$1{-}\varepsilon$};
    \draw (1.0, 0.015) -- (1.0,-0.015) node[above, font=\tiny] {$1$};
    \draw (1.2, 0.015) -- (1.2,-0.015) node[above, font=\tiny] {$1{+}\varepsilon$};

    % Вертикальная пунктирная при 1-eps
    \draw[dashed, gray!60] (0.8, 0) -- (0.8, -0.325);

    \node[font=\small, align=center] at (1.25, -0.95)
          {$\hat{A}_t < 0$: отсечение при $1-\varepsilon$};
  \end{scope}
\end{tikzpicture}
\caption{Механизм отсечения PPO. \emph{Слева}: при $\hat{A}_t > 0$ целевая функция ограничена сверху при $r_t = 1+\varepsilon$ (синий), предотвращая чрезмерный рост вероятности действия. \emph{Справа}: при $\hat{A}_t < 0$ целевая функция ограничена снизу при $r_t = 1-\varepsilon$ (красный), предотвращая чрезмерное снижение. В обоих случаях оператор $\min$ делает $L^{\mathrm{CLIP}}$ пессимистической нижней границей необрезанного суррогата (серый пунктир).}
\label{fig:ppo-clip}
\end{figure}

\subsection{Полная целевая функция PPO}

На практике PPO оптимизирует составную целевую функцию, сочетающую отсечённые потери актора с потерями регрессии критика и бонусом за энтропию:
\begin{equation}
\label{eq:ppo-full}
\mathcal{L}(\theta, \phi)
= -L^{\mathrm{CLIP}}(\theta)
  + c_v\,\E\bigl[(V_\phi(S_t) - \hat{V}_t)^2\bigr]
  - c_H\,\E\bigl[-\log\pi_\theta(A_t \mid S_t)\bigr],
\end{equation}
где $\hat{V}_t$ \,---\, цель GAE~\eqref{eq:gae-vtarget}, $c_v > 0$ взвешивает потери ценности, а $c_H > 0$ взвешивает бонус за энтропию $H(\pi_\theta({\cdot}|S_t))$. Три члена выполняют различные функции:
\begin{description}[leftmargin=!, labelwidth=3.5cm]
  \item[$-L^{\mathrm{CLIP}}$] Максимизирует целевую функцию стратегии с консервативным ограничением доверительной области.
  \item[$c_v\,L^{\mathrm{VF}}$] Обучает критика предсказывать $V^{\pi_\theta}$, что необходимо для точной оценки преимущества.
  \item[$-c_H\,H$] Добавляет бонус за энтропию, препятствующий преждевременному коллапсу стратегии к почти детерминированному распределению, поддерживая исследование в ходе обучения.
\end{description}
Когда актор и критик разделяют слои сети, вся целевая функция~\eqref{eq:ppo-full} дифференцируется совместно.

\subsection{Алгоритм PPO}

\begin{algorithm}[ht]
\caption{Proximal Policy Optimization (PPO)}
\label{alg:ppo}
\DontPrintSemicolon
\KwIn{Длина прокрутки $T$, рабочие $K$, эпохи $K_e$, размер мини-пакета $M$, отсечение $\varepsilon$, $\lambda$ для GAE}
Инициализировать $\theta$, $\phi$\;
\Repeat{сходимость}{
  \tcp{Сбор данных}
  \For{каждый из $K$ рабочих параллельно}{
    Запустить $\pi_\theta$ на $T$ шагах; сохранить $\{S_t, A_t, R_{t+1}, S_{t+1}\}$\;
    Вычислить $\hat{A}_t^{\mathrm{GAE}}$ через~\eqref{eq:gae-def} и $\hat{V}_t$ через~\eqref{eq:gae-vtarget}\;
  }
  Нормализовать преимущества: $\hat{A}_t \leftarrow (\hat{A}_t - \bar{A})/\hat{\sigma}_A$ по всем рабочим\;
  Зафиксировать $\theta_\mathrm{old} \leftarrow \theta$\;
  \tcp{Несколько эпох оптимизации по собранным данным}
  \For{эпоха $= 1, \ldots, K_e$}{
    \For{каждый мини-пакет $\mathcal{B}$ размера $M$}{
      Вычислить $r_t(\theta) = \pi_\theta(A_t|S_t)/\pi_{\theta_\mathrm{old}}(A_t|S_t)$ для каждого $(S_t, A_t) \in \mathcal{B}$\;
      Вычислить $\mathcal{L}(\theta,\phi)$ через~\eqref{eq:ppo-full}\;
      Обновить $\theta, \phi$ одним шагом Adam\;
    }
  }
}
\end{algorithm}

Несколько проектных решений в алгоритме~\ref{alg:ppo} заслуживают комментария.

\begin{remark}[Нормализация преимущества]
Вычитание среднего по пакету и деление на стандартное отклонение по пакету (шаг «Нормализовать преимущества») является стандартной инженерной практикой. Это поддерживает примерно постоянную величину градиента на протяжении эпизодов независимо от масштаба вознаграждений, снижая чувствительность к скорости обучения и стабилизируя обучение. Это вносит небольшое смещение, поскольку нормализация меняет относительные веса переходов, но эмпирически полезна.
\end{remark}

\begin{remark}[Несколько эпох]
PPO повторно использует каждый пакет для $K_e$ шагов градиента. Это существенно улучшает эффективность использования выборок по сравнению с чистыми методами обучения на политике (где данные отбрасываются после одного шага градиента), тогда как отсечение гарантирует, что стратегия не отклоняется слишком далеко от стратегии сбора данных в рамках одного пакета. Типичные значения: $K_e \in [3, 10]$.
\end{remark}

\begin{example}[PPO на простом МПР]
\label{ex:ppo-numeric}
Рассмотрим МПР в виде цепочки из трёх состояний $\{s_1, s_2, s_3\}$, где $s_3$ является терминальным. У агента есть два действия (влево/вправо). Предположим, на текущей итерации:
\begin{itemize}[leftmargin=*]
  \item $\hat{A}_t = 0.8$ (совершённое действие было значительно выше среднего).
  \item $r_t(\theta) = 1.35$ после шага градиента (новая стратегия назначает на $35\%$ больше вероятности этому действию, чем старая).
  \item $\varepsilon = 0.2$.
\end{itemize}
Без отсечения вклад в $L$ составил бы $1.35 \times 0.8 = 1.08$. С отсечением $r_t$ ограничивается до $1+\varepsilon = 1.2$, давая $1.2 \times 0.8 = 0.96$. Поскольку $0.96 < 1.08$, оператор $\min$ выбирает $0.96$. Градиент $L^{\mathrm{CLIP}}$ по $\theta$ равен нулю для этой выборки, потому что $r_t > 1+\varepsilon$: стратегия уже достаточно сдвинулась в направлении увеличения вероятности этого действия, и PPO не даёт дополнительного градиента для дальнейшего смещения в этом направлении.\qed
\end{example}

% -----------------------------------------------------------
\section{PPO для языковых моделей}
\label{sec:ppo-llm}
% -----------------------------------------------------------

\subsection{Языковые модели как стратегии ОП}

Применение PPO для дообучения больших языковых моделей (LLM) с использованием обратной связи от людей\,---\,обычно называемое RLHF\index{RLHF}\index{Reinforcement Learning from Human Feedback|see{RLHF}}\,---\,стало одним из наиболее практически важных применений обучения с подкреплением. Глава~\ref{ch:lm-post-training} подробно изучает языковые модели как агентов ОП; здесь мы сосредотачиваемся конкретно на компоненте PPO в пайплайне RLHF.

Языковая модель $\pi_\theta$ определяет авторегрессионную стратегию по токенам. При промпте $x = (x_1, \ldots, x_m)$ модель генерирует ответ $y = (y_1, \ldots, y_T)$ по одному токену за раз:
\begin{equation}
\label{eq:lm-policy}
\pi_\theta(y \mid x)
= \prod_{t=1}^T \pi_\theta(y_t \mid x,\, y_1, \ldots, y_{t-1}).
\end{equation}
Это соответствует МПР, где состоянием на шаге $t$ является $S_t = (x, y_1, \ldots, y_{t-1})$, действием является следующий токен $A_t = y_t \in \mathcal{V}$ (словарь), а эпизод завершается при генерации токена конца последовательности. Пространство действий имеет размер $|\mathcal{V}| \in [32{,}000,\; 128{,}000]$\,---\,значительно больше, чем в классических задачах ОП,\,---\,что делает методы на основе функции ценности неприемлемыми, а стратегические методы, такие как PPO, естественными.

\subsection{Моделирование вознаграждения и штраф KL}

В RLHF сигнал вознаграждения поступает от отдельно обученной \textbf{модели вознаграждения}\index{reward model} $r_\psi(x, y)$, обученной на сравнениях предпочтений людей. В простейшем случае скалярное вознаграждение выдаётся только в конце эпизода: $R_t = 0$ при $t < T$ и $R_T = r_\psi(x, y)$.

Однако без дополнительной регуляризации стратегия быстро научается эксплуатировать модель вознаграждения\,---\,явление, называемое \textbf{взломом вознаграждения}\index{reward hacking}. Стратегия может производить выходные данные, которые высоко оцениваются по $r_\psi$, но расходятся от связного, безопасного и фактически точного языка. Стандартным средством является пошаговый штраф за KL-дивергенцию относительно \emph{референсной стратегии}\index{reference policy} $\pi_{\mathrm{ref}}$ (обычно модели после обучения с учителем до обучения ОП):
\begin{equation}
\label{eq:rlhf-reward}
\tilde{R}_t(x, y)
=
\begin{cases}
r_\psi(x, y)
  - \beta\,\log\dfrac{\pi_\theta(y_T \mid x, y_{<T})}
                      {\pi_{\mathrm{ref}}(y_T \mid x, y_{<T})},
& t = T, \\[8pt]
- \beta\,\log\dfrac{\pi_\theta(y_t \mid x, y_{<t})}
                    {\pi_{\mathrm{ref}}(y_t \mid x, y_{<t})},
& t < T,
\end{cases}
\end{equation}
где $\beta > 0$ управляет силой штрафа KL.

\begin{proposition}[Суммарное вознаграждение RLHF]
\label{prop:rlhf-kl}
Сумма скорректированных пошаговых вознаграждений удовлетворяет
\begin{equation}
\label{eq:rlhf-total}
\sum_{t=1}^T \tilde{R}_t
= r_\psi(x, y)
  - \beta\,\KL{\pi_\theta(\cdot \mid x)}{\pi_{\mathrm{ref}}(\cdot \mid x)},
\end{equation}
где KL-дивергенция вычисляется по совместному последовательному распределению.
\end{proposition}

\begin{proof}
Суммируя~\eqref{eq:rlhf-reward} по $t = 1, \ldots, T$:
\[
\sum_{t=1}^T \tilde{R}_t
= r_\psi(x,y)
  - \beta\sum_{t=1}^T
    \log\frac{\pi_\theta(y_t \mid x, y_{<t})}
             {\pi_{\mathrm{ref}}(y_t \mid x, y_{<t})}.
\]
KL-дивергенция между двумя авторегрессионными распределениями $\pi_\theta(\cdot|x)$ и $\pi_{\mathrm{ref}}(\cdot|x)$ факторизуется как
\[
\KL{\pi_\theta(\cdot|x)}{\pi_{\mathrm{ref}}(\cdot|x)}
= \E_{y \sim \pi_\theta}\!\left[
  \sum_{t=1}^T \log\frac{\pi_\theta(y_t|x,y_{<t})}
                         {\pi_{\mathrm{ref}}(y_t|x,y_{<t})}
\right],
\]
поэтому пошаговая выборочная сумма $\sum_t \log[\pi_\theta/\pi_{\mathrm{ref}}]$ является несмещённой оценкой KL методом Монте-Карло. Подстановка даёт~\eqref{eq:rlhf-total}.
\end{proof}

Утверждение~\ref{prop:rlhf-kl} раскрывает целевую функцию, которую RLHF неявно оптимизирует: максимизировать оценку модели вознаграждения минус взвешенный KL-штраф за отклонение от референсной стратегии. При $\beta \to 0$ штраф исчезает, и стратегия может произвольно отклоняться от $\pi_{\mathrm{ref}}$, рискуя взломом вознаграждения. При $\beta \to \infty$ оптимальной стратегией является $\pi_{\mathrm{ref}}$, и обучение ОП не происходит.

\subsection{Четырёхмодельная архитектура RLHF}

Стандартный цикл обучения PPO-RLHF одновременно поддерживает четыре нейронные сети:

\begin{description}[leftmargin=!, labelwidth=3.5cm]
  \item[Актор $\pi_\theta$] Обучаемая стратегия. Её параметры обновляются на каждой итерации PPO.
  \item[Критик $V_\phi$] Голова функции ценности\,---\,обычно один линейный слой, применённый к скрытому состоянию последнего токена\,---\,оценивающая ожидаемую накопленную скорректированную отдачу из состояния $S_t$.
  \item[Референсная модель $\pi_{\mathrm{ref}}$] Замороженная копия модели после обучения с учителем, используемая только для вычисления штрафа KL в~\eqref{eq:rlhf-reward}.
  \item[Модель вознаграждения $r_\psi$] Замороженная модель (обычно той же архитектуры, что и актор, с добавленной скалярной головой), обученная на данных предпочтений людей для выдачи скалярной оценки качества.
\end{description}

Преимущества $\hat{A}_t$ вычисляются через GAE, применённый к скорректированным вознаграждениям $\tilde{R}_t$ из~\eqref{eq:rlhf-reward}, а отсечённая целевая функция PPO~\eqref{eq:ppo-clip} применяется с отношениями важности на уровне словаря $r_t(\theta) = \pi_\theta(y_t|x,y_{<t}) / \pi_{\theta_\mathrm{old}}(y_t|x,y_{<t})$.

\begin{example}[PPO-RLHF на задаче ответа на вопросы]
\label{ex:ppo-rlhf}
Предположим, промпт \,---\, «Объясни последнюю теорему Ферма.», а стратегия генерирует ответ из 120 токенов $y$. Модель вознаграждения оценивает полный ответ: $r_\psi(x,y) = 0.78$. На каждом шаге $t < 120$ вычисляется пошаговое логарифмическое отношение $\log[\pi_\theta(y_t|x,y_{<t})/\pi_{\mathrm{ref}}(y_t|x,y_{<t})]$; если стратегия умеренно отклонилась от референсной, каждый член может составлять около $0.01$, вкладывая суммарный штраф KL $\beta \times 0.01 \times 120 = 1.2\beta$. При $\beta = 0.05$ суммарное скорректированное вознаграждение равно $0.78 - 0.06 = 0.72$.

Критик $V_\phi(S_t)$ предсказывает ожидаемую суммарную скорректированную отдачу из частичной последовательности, видимой до сих пор. После генерации каждого токена вычисляется $\delta_t = \tilde{R}_t + \gamma V_\phi(S_{t+1}) - V_\phi(S_t)$, и GAE агрегирует эти ошибки в пошаговые оценки преимущества. Токен вроде «Уайлс» (имея в виду Эндрю Уайлса), который одновременно уместен и специфичен, может повысить предсказание критика, давая положительное преимущество, повышающее вероятность этого токена в будущих обновлениях PPO.\qed
\end{example}

\subsection{Практические соображения}

Следующие рекомендации дистиллируют эмпирические лучшие практики из литературы по RLHF (Оуян и~др.\ 2022, Бай и~др.\ 2022, Циглер и~др.\ 2019).

\paragraph{Мониторинг и управление KL.}
Отслеживайте скользящее среднее $\KL{\pi_\theta}{\pi_{\mathrm{ref}}}$ по пакетам промптов. Быстрый рост сигнализирует об отклонении стратегии от референсной модели. Если KL превышает предустановленный порог (например, 10 натов), либо увеличьте $\beta$, уменьшите скорость обучения, либо остановите обучение и диагностируйте модель вознаграждения.

\paragraph{Нормализация преимущества.}
Нормализуйте $\hat{A}_t$ до нулевого среднего и единичной дисперсии по мини-пакету перед вычислением отсечённой целевой функции. Для генерации языка эпизоды (ответы) имеют различную длину; нормализация предотвращает доминирование длинных последовательностей в обновлениях градиента исключительно из-за накопления большего числа пошаговых членов.

\paragraph{Инициализация критика.}
Инициализируйте голову функции ценности поверх того же предобученного каркаса, что и актор. Запуск со случайно инициализированным критиком даёт дико неточные оценки преимущества в начале обучения, дестабилизируя актора до того, как он что-либо усвоит.

\paragraph{Количество эпох на пакет.}
Используйте $K_e \in [1, 4]$ эпох оптимизации на пакет данных для RLHF против $K_e \in [4, 10]$ для классических сред ОП. Поскольку модель вознаграждения может эксплуатироваться (в отличие от настоящей среды), слишком много эпох на одном пакете побуждает актора переобучаться к сигналу модели вознаграждения, ускоряя взлом вознаграждения.

\paragraph{Обнаружение взлома вознаграждения.}
Периодически оценивайте обученную стратегию на отложенных промптах, оцениваемых людьми (или отдельным независимым оценщиком). Если оценка модели вознаграждения продолжает расти, а оценки людей выравниваются или снижаются, происходит взлом вознаграждения. Меры противодействия включают обновление модели вознаграждения с дополнительными данными предпочтений, увеличение $\beta$ или включение потерь обучения с учителем наряду с целевой функцией PPO.

% -----------------------------------------------------------
\section{Сходимость и практические соображения}
\label{sec:ac-convergence}
% -----------------------------------------------------------

\subsection{Свойства сходимости}

Теория сходимости для методов актор-критик существенно сложнее, чем для чистых методов TD или чистых методов градиента стратегии, поскольку оба компонента обновляются одновременно и их целевые функции взаимодействуют.

В случае \emph{линейной} аппроксимации функции при фиксированной стратегии критик TD(0) сходится к решению проецированного уравнения Беллмана, как установлено в главе~\ref{ch:td}. Когда актор также обновляется, сходимость требует \textbf{двухшкальной стохастической аппроксимации}: размер шага критика $\alpha_\phi$ должен убывать быстрее, чем шаг актора $\alpha_\theta$, обеспечивая, что критик отслеживает $V^{\pi_\theta}$ быстрее, чем меняется $\pi_\theta$.

\begin{theorem}[Сходимость при двух шкалах, набросок]
\label{thm:two-timescale}
\index{actor-critic!two-timescale convergence}
При следующих условиях:
\begin{enumerate}[label=(\roman*), leftmargin=*]
  \item линейная аппроксимация функции для критика,
  \item размеры шагов Роббинса-Монро, удовлетворяющие $\sum_t \alpha_\phi(t) = \infty$, $\sum_t \alpha_\phi(t)^2 < \infty$, и $\alpha_\theta(t) / \alpha_\phi(t) \to 0$,
  \item обыкновенное дифференциальное уравнение градиента стратегии имеет единственное глобально притягивающее равновесие (или, в более общем случае, итерации остаются ограниченными),
\end{enumerate}
алгоритм актор-критик сходится п.н.\ к стационарной точке целевой функции $J(\theta)$.
\end{theorem}

Для \emph{нелинейной} аппроксимации функции (нейронные сети) сопоставимо чёткой теории не существует. На практике актор и критик обновляются одним оптимизатором (например, Adam) с одинаковыми или схожими скоростями обучения, что нарушает требование двух шкал. Тем не менее PPO эмпирически стабилен на широком классе задач, что предполагает, что нейронные сети действуют как собственные эффективные регуляризаторы.

\subsection{Сравнение алгоритмов градиента стратегии}

В таблице~\ref{tab:pg-comparison} обобщены ключевые свойства семейства методов градиента стратегии, обсуждавшихся в главах~\ref{ch:pg} и~\ref{ch:actor-critic}.

\begin{table}[ht]
\centering
\caption{Сравнение алгоритмов градиента стратегии. «Обучение на политике» означает, что каждая выборка данных используется не более чем для одного обновления градиента.}
\label{tab:pg-comparison}
\footnotesize
\begin{tabular}{@{}lp{2.2cm}ccc@{}}
\toprule
Алгоритм & Базовая линия & Онлайн? & Неск.\ эпох? & Довер.\ обл.? \\
\midrule
REINFORCE               & Нет / конст.  & Нет  & Нет  & Нет  \\
REINFORCE + баз. линия & $V_\phi(s)$      & Нет  & Нет  & Нет  \\
Одношаговый АК          & Критик TD(0)     & Да & Нет  & Нет  \\
A2C / A3C               & Критик GAE       & Да & Нет  & Нет  \\
TRPO                    & Критик GAE       & Да & Нет & KL \\
PPO                     & Критик GAE       & Да & Да & clip~$\varepsilon$ \\
\bottomrule
\end{tabular}
\end{table}

\subsection{Чувствительность к гиперпараметрам и советы по реализации}

PPO заметно устойчив к выбору гиперпараметров по сравнению с TRPO, который требует тщательной настройки числа шагов сопряжённых градиентов и параметров поиска по направлению. Тем не менее следующие рекомендации полезны на практике.

\begin{description}[leftmargin=!, labelwidth=3.5cm]
  \item[Параметр отсечения $\varepsilon$] Типичные значения: от $0.1$ до $0.3$. Меньшее $\varepsilon$ даёт более консервативные обновления и безопаснее при дообучении крупных предобученных моделей. $\varepsilon = 0.2$ \,---\, наиболее распространённое значение по умолчанию.
  \item[GAE $\lambda$] Значения $0.90$--$0.97$ хорошо работают для большинства задач. Более длинные эпизоды требуют большего $\lambda$ (для захвата долгосрочного кредита).
  \item[Коэффициент ценности $c_v$] Обычно $0.5$; можно уменьшить, если сеть функции ценности оптимизируется отдельно или если величины градиентов уже хорошо согласованы.
  \item[Коэффициент энтропии $c_H$] Начинать с $0.01$; увеличивать, если стратегия преждевременно коллапсирует к почти детерминированному распределению.
  \item[Эпохи на пакет $K_e$] Значения $4$--$8$ для непрерывного управления; $1$--$4$ для RLHF. Мониторить KL между $\pi_\theta$ и $\pi_{\theta_\mathrm{old}}$ и останавливать эпохи досрочно, если KL превышает $\sim 1.5\varepsilon$.
  \item[Размер пакета и мини-пакета] Большие пакеты снижают шум градиента, но требуют больше памяти. Прокрутка из $2048$ шагов с размером мини-пакета $64$ \,---\, распространённая отправная точка для непрерывного управления.
  \item[Скорость обучения] $3 \times 10^{-4}$ с Adam \,---\, устойчивое значение по умолчанию. Линейные или косинусные расписания затухания часто помогают при дообучении.
  \item[Нормализация наблюдений] Нормализовать наблюдения до нулевого среднего и единичной дисперсии (с накопленной статистикой) и ограничивать вознаграждения фиксированным диапазоном (например, $[-5, 5]$), чтобы критик не был перегружен выбросами вознаграждений.
\end{description}

% -----------------------------------------------------------
\section*{Упражнения}
\addcontentsline{toc}{section}{Упражнения}
\label{sec:ch9-exercises}
% -----------------------------------------------------------

\begin{exercise}
\label{ex:advantage-computation}
Рассмотрим МПР с тремя действиями, доступными в состоянии $s$. Предположим, $Q^{\pi}(s, a_1) = 5$, $Q^{\pi}(s, a_2) = 8$, $Q^{\pi}(s, a_3) = 2$, и стратегия назначает вероятности $\pi(a_1|s) = 0.5$, $\pi(a_2|s) = 0.3$, $\pi(a_3|s) = 0.2$.
\begin{enumerate}[label=(\alph*)]
  \item Вычислите $V^\pi(s)$ и преимущество $A^\pi(s, a_i)$ для каждого действия.
  \item Проверьте, что $\E_{A \sim \pi(\cdot|s)}[A^\pi(s,A)] = 0$.
  \item Шаг градиента увеличивает $\pi(a_2|s)$. Объясните, почему использование $A^\pi(s,a_2) = 2$ (вместо $Q^\pi(s,a_2) = 8$) в качестве веса градиента приводит к меньшему, но более информативному обновлению.
\end{enumerate}
\end{exercise}

\begin{exercise}
\label{ex:td-advantage-numeric}
Критик с $V_\phi = V^{\pi_\theta}$ наблюдает переход $(S_t, A_t, R_{t+1}, S_{t+1})$ с $R_{t+1} = 2$, $\gamma = 0.9$, $V^\pi(S_t) = 3$, $V^\pi(S_{t+1}) = 4$.
\begin{enumerate}[label=(\alph*)]
  \item Вычислите ошибку TD $\delta_t$.
  \item Действие $A_t$ выше или ниже среднего в состоянии $S_t$? Обоснуйте, используя знак $\delta_t$.
  \item После одного обновления актора с размером шага $\alpha_\theta = 0.01$ и $\gamma^t = 1$, на сколько изменится $\log \pi_\theta(A_t|S_t)$?
\end{enumerate}
\end{exercise}

\begin{exercise}
\label{ex:gae-derivation}
Пусть $T = 4$ (шаги $t = 0,1,2,3$, терминальное при $t=4$) с ошибками TD $\delta_0 = 0.4$, $\delta_1 = -0.2$, $\delta_2 = 0.6$, $\delta_3 = 0.1$. Возьмём $\gamma = 1$ и $\lambda = 0.8$.
\begin{enumerate}[label=(\alph*)]
  \item Вычислите $\hat{A}_0^{\mathrm{GAE}}$ используя определение~\ref{def:gae-ch9}.
  \item Проверьте, что ваш ответ совпадает с формулой из утверждения~\ref{prop:gae-nstep}, вычислив взвешенную смесь $n$-шаговых преимуществ $\hat{A}_0^{(n)}$ для $n = 1, 2, 3, 4$.
  \item Вычислите $\hat{A}_0^{\mathrm{GAE}}$ для $\lambda = 0$ и $\lambda = 1$ и интерпретируйте результаты.
\end{enumerate}
\end{exercise}

\begin{exercise}
\label{ex:ppo-clip-numerical}
Рассмотрим отсечённую целевую функцию PPO с $\varepsilon = 0.2$.
\begin{enumerate}[label=(\alph*)]
  \item Вычислите вклады $L^{\mathrm{CLIP}}$ для следующих трёх переходов и укажите, активно ли отсечение:
    \begin{enumerate}[label=(\roman*)]
      \item $\hat{A}_t = 1.5$, $r_t(\theta) = 0.8$.
      \item $\hat{A}_t = 1.5$, $r_t(\theta) = 1.3$.
      \item $\hat{A}_t = -0.5$, $r_t(\theta) = 0.75$.
    \end{enumerate}
  \item Для случая (ii) объясните, почему градиент $L^{\mathrm{CLIP}}$ по $\theta$ равен нулю, несмотря на то что $\hat{A}_t > 0$ и действие следует поощрять. Почему такое консервативное поведение желательно?
\end{enumerate}
\end{exercise}

\begin{exercise}
\label{ex:surrogate-gradient-proof}
В этом упражнении детально разбирается суррогатный градиент (утверждение~\ref{prop:surrogate-gradient}).
\begin{enumerate}[label=(\alph*)]
  \item Покажите, что $r_t(\theta)|_{\theta=\theta_\mathrm{old}} = 1$ для всех переходов.
  \item Запишите $\nabla_\theta r_t(\theta)$ через стратегию и покажите, что оно равно $r_t(\theta)\,\nabla_\theta \log\pi_\theta(A_t|S_t)$.
  \item Заключите, что $\nabla_\theta L(\theta)|_{\theta_\mathrm{old}} = \E[\nabla_\theta\log\pi_{\theta_\mathrm{old}}(A_t|S_t)\,\hat{A}_t]$, и объясните, почему это в точности градиент стратегии.
\end{enumerate}
\end{exercise}

\begin{exercise}
\label{ex:rlhf-kl}
В скорректированном вознаграждении RLHF~\eqref{eq:rlhf-reward} задать $\gamma = 1$.
\begin{enumerate}[label=(\alph*)]
  \item Докажите утверждение~\ref{prop:rlhf-kl} с нуля, используя только цепное правило для KL-дивергенции авторегрессионных моделей.
  \item Покажите, что оптимальная стратегия для целевой функции $\max_\pi \E[r_\psi(x,y)] - \beta\,\KL{\pi(\cdot|x)}{\pi_{\mathrm{ref}}(\cdot|x)}$ есть $\pi^*(y|x) \propto \pi_{\mathrm{ref}}(y|x)\,e^{r_\psi(x,y)/\beta}$.
  \item Интерпретируйте оптимальную стратегию: что происходит при $\beta \to 0$? При $\beta \to \infty$? Что это означает для выбора $\beta$ на практике?
\end{enumerate}
\end{exercise}

\begin{exercise}
\label{ex:ac-comparison-table}
Заполните следующую таблицу сравнения для трёх алгоритмов: REINFORCE с базовой линией, одношаговый актор-критик и PPO. Для каждой строки ответьте «да» / «нет» или дайте краткую фразу.

\begin{center}
\begin{tabular}{@{}lccc@{}}
\toprule
Свойство & REINFORCE+БЛ & Одношаг. АК & PPO \\
\midrule
Требует полного эпизода?         & & & \\
Повторное использование данных?  & & & \\
Оценщик преимущества             & & & \\
Обновления онлайн (пошагово)?    & & & \\
Ограничение доверительной обл.?  & & & \\
\bottomrule
\end{tabular}
\end{center}
\end{exercise}

\begin{exercise}
\label{ex:gae-bias-variance}
Объясните, используя разложение из утверждения~\ref{prop:gae-nstep}, почему GAE с $\lambda < 1$ имеет меньшую дисперсию, чем оценка преимущества Монте-Карло $G_t - V_\phi(S_t)$, но вносит смещение.
\begin{enumerate}[label=(\alph*)]
  \item Определите источник дисперсии в оценке Монте-Карло (подсказка: подсчитайте количество независимых случайных величин в $G_t$).
  \item Определите источник смещения в оценке GAE при $\lambda < 1$ (подсказка: рассмотрите, что происходит, когда $V_\phi \neq V^\pi$).
  \item Неформально обоснуйте, почему $\lambda \approx 0.95$ является разумным значением по умолчанию для эпизодов длиной $\sim 1000$.
\end{enumerate}
\end{exercise}

\begin{exercise}
\label{ex:two-timescale}
Результат сходимости при двух шкалах (теорема~\ref{thm:two-timescale}) требует $\alpha_\theta(t)/\alpha_\phi(t) \to 0$.
\begin{enumerate}[label=(\alph*)]
  \item Объясните интуитивно, почему критик должен обучаться быстрее актора. Что пойдёт не так, если обновления актора такие же быстрые, как у критика?
  \item Предложите конкретную пару последовательностей размеров шагов $\alpha_\phi(t) = 1/(t+1)$ и $\alpha_\theta(t)$, удовлетворяющих условию двух шкал.
  \item В современных реализациях PPO с Adam и актор, и критик используют одинаковый оптимизатор и скорость обучения. Какой практический механизм заменяет требование двух шкал?
\end{enumerate}
\end{exercise}

\begin{exercise}[Повышенная сложность]
\label{ex:trpo-ppo-connection}
В этом упражнении устанавливается точная связь между TRPO и PPO.
\begin{enumerate}[label=(\alph*)]
  \item Покажите, что ограничение TRPO $\KL{\pi_\theta}{\pi_{\theta_\mathrm{old}}} \leq \delta$ влечёт $|r_t(\theta) - 1| \leq \sqrt{2\delta}$ для любой отдельной пары действие-состояние, используя неравенство Пинскера $\|\pi - q\|_1 \leq \sqrt{2\,\KL{\pi}{q}}$.
  \item Используя часть (a), покажите, что отсечение PPO при $\varepsilon$ является приближением первого порядка к ограничению KL TRPO с $\delta \approx \varepsilon^2/2$.
  \item Обсудите одно преимущество и один недостаток приближения PPO по сравнению с точным TRPO.
\end{enumerate}
\end{exercise}
